
\begin{table}[!h]
   \centering
   \small
   \begin{threeparttable}[b]
      \caption{\label{tab:encaje_loan_growth_lender_fe} Reserve-requirement shocks and credit growth with lender fixed effects}
      \bigskip
      \begin{tabular}{lcccc}
         \toprule
         Credit type & \multicolumn{2}{c}{Credit lines} & \multicolumn{2}{c}{Term loans} \\ 
         \cmidrule(lr){2-3} \cmidrule(lr){4-5}
         Outcome & \multicolumn{2}{c}{Commitment} & \multicolumn{2}{c}{Drawn} \\ 
                                           & (1)                   & (2)                   & (3)                   & (4)\\  
         \midrule 
         Bank shock                        & 0.593$^{*}$           & 0.349$^{**}$          & -0.028                & -0.064$^{**}$\\   
                                           & (0.290)               & (0.143)               & (0.029)               & (0.027)\\   
         Lenders (= clusters)              & 16                    & 16                    & 35                    & 35\\  
         Firms                             & 2,913                 & 2,913                 & 3,632                 & 3,632\\  
         Firm-by-date cells                & 131,729               & 131,729               & 137,974               & 137,974\\  
         Wild cluster bootstrap $p$-value  & 0.143                 & 0.136                 & 0.399                 & 0.078\\  
         Mean of dependent variable        & 0.019                 & 0.019                 & -0.045                & -0.045\\  
          \\
         Observations                      & 305,776               & 305,776               & 332,487               & 332,487\\  
         R$^2$                             & 0.407                 & 0.409                 & 0.410                 & 0.411\\  
         Within R$^2$                      & $5.56\times 10^{-5}$  & $1.86\times 10^{-5}$  & $2.53\times 10^{-6}$  & $1.29\times 10^{-5}$\\   
          \\
         Firm $\times$ date fixed effects  & $\checkmark$          & $\checkmark$          & $\checkmark$          & $\checkmark$\\   
         Lender fixed effects              &                       & $\checkmark$          &                       & $\checkmark$\\   
         \bottomrule
      \end{tabular}
      
      \begin{tablenotes}\item \textit{Notes:} The dependent variable is the log change between consecutive month-ends of a firm's position at a lender. For credit lines (columns 1--2) it is the total commitment, drawn credit (\textit{cr\'edito vigente}) plus the unused line (\textit{l\'inea no utilizada}); for term loans (columns 3--4) it is drawn credit. The shock is the minimum legal reserve rate shock. Columns (1) and (3) reproduce the corresponding columns of the main table; columns (2) and (4) add lender fixed effects, so the coefficient is identified only from changes within a lender over time, net of its average growth. Lenders are all institutions with a reserve shock: commercial banks, cajas municipales and rurales, and financieras. Credit is classified with a manual product table: credit cards, current-account advances, guarantees, letters of credit, credits to be settled, multi-line facilities and other lines are credit lines; leasing, lease-back, real-estate, mortgage, pawn, income-property and RFA credit are term loans; \textit{pr\'estamos} are credit lines when reported as revolving, term loans when non-revolving, and, when the revolving status is not identified, credit lines if the row carries an unused line and term loans otherwise. Products whose nature is ambiguous (overdrafts, factoring, discounting, foreign-trade finance, project finance, other credits, credit to consolidable entities) are excluded from all columns. Amounts are summed across products and currencies within a bank--firm--credit-type cell (foreign-currency amounts are reported in soles at the month-end exchange rate, so growth includes exchange-rate revaluation). Growth is computed for continuing relationships with strictly positive amounts at both dates; entry and exit are not captured. The bank shock is $s^{mn}_{b,0}\Delta r^{mn}_{t}+s^{me}_{b,0}\Delta r^{me}_{t}$, where $\Delta r$ is the month-end change, in percentage points, of the statutory minimum legal reserve rate or marginal reserve rate on domestic (mn) and foreign (me) currency obligations (positive = tighter), and $s_{b,0}$ are the bank's fixed December 2009 deposit shares in each currency. The marginal-rate shock is undefined when the domestic regime is a single rate, which shortens the marginal-rate columns. Each regression keeps firm-by-date cells with at least two lenders and absorbs firm-by-date fixed effects. Sample: RCC sample of 10,000 firms, September 2010 to December 2023. Standard errors, in parentheses, are clustered by lender; with few clusters the $p$-value from a wild cluster bootstrap (null imposed, 9,999 draws) is reported. * $p<0.10$, ** $p<0.05$, *** $p<0.01$ from the conventional clustered $t(G-1)$ test.
      \end{tablenotes}
   \end{threeparttable}
\end{table}


